\documentclass[10pt,a4paper]{amsart}
\usepackage[margin=19mm]{geometry}
\usepackage{amsmath,amssymb,mathtools}
\usepackage[scr=rsfs,cal=euler]{mathalfa}
\usepackage{xfrac}
\usepackage[unicode,hidelinks,bookmarks=false]{hyperref}
\newcommand{\ie}{i.e.\ }
\newcommand{\cf}{cf.\ }
\newcommand{\resp}{respectively\ }
\newcommand{\Subsection}{\subsection}
\providecommand{\nobreakdash}{\nobreak\hbox{-}}
\setlength{\emergencystretch}{2em}
\multlinegap0pt
\usepackage{bm}
\usepackage{bbm}
\usepackage{dsfont}
\usepackage{xspace}
\usepackage{cancel}
\usepackage{mathtools}
\usepackage{cleveref}
\usepackage[shortlabels]{enumitem}
\usepackage{amsthm}
\makeatletter
\@ifundefined{theorem}{\newtheorem{theorem}{Theorem}}{}
\@ifundefined{lemma}{\newtheorem{lemma}[theorem]{Lemma}}{}
\makeatother
\makeatletter
\renewcommand{\a}{\alpha}
\renewcommand{\b}{\beta}
\renewcommand{\binom}[2]{\left(\begin{smallmatrix}#1\\\\#2\end{smallmatrix}\right)}
\renewcommand{\d}{{\delta}}
\expandafter\ifx\csname discussion\endcsname\relax
\newenvironment{discussion}{\ifdefs\noindent \newline \noindent \bf
	************************** begin discussion\rm \\ \bf\fi}{\ifdefs \rm \noindent \newline
	\noindent \bf ************************** end discussion \rm\\ \fi}
\fi
\expandafter\ifx\csname extradetails\endcsname\relax
\newenvironment{extradetails}{\ifdefs \noindent \newline \noindent \bf
	************************** begin extra details\rm \\ \bf\fi }{\ifdefs \rm \noindent \newline
	\noindent \bf ************************** end extra details \rm\\ \fi}
\fi
\renewcommand{\k}{\kappa}
\renewcommand{\l}{\lambda}
\newcommand{\smallbinom}[2]{(\begin{smallmatrix}#1\\#2\end{smallmatrix})}
\makeatother
\title[The Jensen--P\'{o}lya program and inequalities for arithmeti]{The Jensen--P\'{o}lya program and inequalities for arithmetic sequences}
\author{Koustav Banerjee, Kathrin Bringmann,, Larry Rolen}
\date{Source: arXiv:2609.14295v2 (CC BY 4.0)}
\begin{document}
\maketitle

\noindent\emph{Source and licence.} The Jensen--P\'{o}lya program and inequalities for arithmetic sequences. Version arXiv:2609.14295v2 (CC BY 4.0), distributed under the Creative Commons Attribution 4.0 International licence, as recorded in the arXiv licence metadata for this exact version and shown on the abstract page https://arxiv.org/abs/2609.14295v2. Full rights notice: licenses/arxiv-2609.14295-CC-BY-4.0.txt.\par\smallskip

\noindent\textit{Editorial scope.} Counted unit: the Lemma whose source label is \texttt{fact2} in the article (source0.tex lines 885--891, source label \texttt{fact2}), delivered together with the article's own complete original proof of it (source0.tex lines 893--952, one \texttt{proof} environment of 60 lines). No mathematics is written, repaired or re-derived: statement and proof are the article's own text line for line. Required context retained uncounted: GenGORZThm1assump1 (lines 250--263); GenGORZThm1assump2 (lines 252--254); InflogconcaveThmassump1 (lines 252--254); genthmassump2 (lines 252--254); InflogconcaveThmassump2 (lines 258--260); InflogconcaveThmassump3 (lines 258--260); genexample (lines 265--267); genexampleassump2 (lines 269--276); genexampleassump3 (lines 269--276); genexampleassump4 (lines 269--276); Fact1 (lines 876--882). Every definition, display and label of those lines is retained unchanged. Omitted from this package and not redistributed: 1--249, 255--257, 261--264, 268--268, 277--875, 883--884, 892--892, 953--1727. Nothing inside a retained excerpt is omitted or altered: every retained body is delivered exactly as the article's lines are written, so no omission receipt of any kind accompanies this package. Citations: the retained excerpts carry no citation command at all, so no bibliography entry of the article is transcribed and no reference is redistributed. Reference adaptation: none was needed and none was made; every reference inside the delivered unit resolves inside it, so no unresolved reference or citation is produced. Printed slips kept literal: no printed word, symbol, hypothesis or number of the counted statement or of its proof was altered. Layout and class: the article's own class is \texttt{amsart} with packages including a4wide, fullpage, graphicx, amsmath, amssymb, amsfonts, amsthm, enumitem, bm, xcolor, cleveref, ulem, url, color; the wrapper is the lane's pinned \texttt{amsart} editorial wrapper at 10pt on A4, which restates the theorem-like environments the retained text needs and adds the article's own macro definitions that the retained text needs (\texttt{\textbackslash a}, \texttt{\textbackslash b}, \texttt{\textbackslash binom}, \texttt{\textbackslash d}, \texttt{\textbackslash k}, \texttt{\textbackslash l}, \texttt{\textbackslash smallbinom}), with extra wrapper packages (bm, bbm, dsfont, xspace, cancel, mathtools, cleveref, enumitem). The delivered numbering is the unit's own. Assets: no retained span contains a figure environment, a graphics-inclusion command, a PDF-page inclusion, a table or a TikZ picture; no image file is copied into this package, the package declares no asset dependency and no image file is redistributed with it. Licence: version arXiv:2609.14295v2 of this article is distributed under the Creative Commons Attribution 4.0 International licence, as recorded in the arXiv licence metadata for this exact version and shown on the abstract page https://arxiv.org/abs/2609.14295v2. A full rights notice is delivered at licenses/arxiv-2609.14295-CC-BY-4.0.txt. Characters: every retained excerpt is pure ASCII, so the delivered engine, XeLaTeX with its default Latin Modern text font, reports no missing character.
 \begin{enumerate}[label=(C\arabic*)]
  \item \label{GenGORZThm1assump1} For $\a_1,\a_2\in \mathbb{R}$, $\a_1\le j\le a_2$, there exist $\{A(n)\}_{n\ge 0}$, $\{\d(n)\}_{n\ge 0}$ such that
 \begin{equation*}%\label{GenGORZthm1eqn1}
 \hspace{2 cm}\log\left(\frac{a(n+j)}{a(n)}\right)=A(n)j-\d^2(n)j^2+o\left(1\right)\ \ \ \ (\text{as}\ n\to \infty).
 \end{equation*}
  \item \label{GenGORZThm1assump2} We have $\lim_{n\to \infty}\d(n)=0$ and $\d(n)>0$ for $n\!\gg\!1$.
 \item \label{genthmassump2} We have $\lim_{n\to \infty}A(n)=0$.
	\item \label{InflogconcaveThmassump1} For $\a_1,\a_2\in \mathbb{R}$ and  $\a_1\le j\le a_2$, there exist $\{B(n)\}_{n\ge 0}$ and $\{\b(n)\}_{n\ge 0}$ such that  %we have, as $n\to \infty$,
	\begin{equation*}%\label{Inflogconcaveeqn2}
	\log\left(\frac{\d(n+j)}{\d(n)}\right)=B(n)j+\b^2(n)j^2+o\left(1\right)\ \ \ \ (\text{as}\ n\to \infty).
	\end{equation*}
	\item \label{InflogconcaveThmassump2} We have $\lim_{n\to \infty}B(n)= 0$.
	\item \label{InflogconcaveThmassump3} We have $\b(n)=o(\d(n))$ as $n\to \infty$.
\end{enumerate}

\begin{equation}\label{genexample}
a_d(n)\sim c_1(d)\frac{\exp\left(\sum_{\l\in \mathcal{S}}A_{\l}(d)n^{\l}\right)}{n^{c_2(d) }}\ \ \ \ (\text{as}\ n\to \infty)
\end{equation}

\begin{enumerate}[label=(C\arabic*), start=7,  labelwidth=\widthof{(C99)}, % Set labelwidth to accommodate the widest expected label
	leftmargin=\labelwidth+\labelsep, % Ensure left margin respects the new labelwidth
	align=left]
		\item \label{genexampleassump1} For $\l\in \mathcal{S}$, we have $A_{\l}(d)\neq 0$.
	\item \label{genexampleassump2} The set $\mathcal{S}\subset\mathbb{Q}^{+}\cap (0,1)$ is finite.
	\item \label{genexampleassump3} We have $A_{\l^{*}}(d)>0$ with $\l^{*}:=\underset{\l \in \mathcal{S}}{\max}\{\l \}$.
	\item \label{genexampleassump4} We have $c_1(d)>0$ and $c_2(d)\in \mathbb{R}$.
\end{enumerate} 

\begin{lemma}\label{Fact1}
	If $f(n)\sim g(n)$ as $n\to \infty$, then, for $\a_1\le j\le \a_2$ %$\mathbb{Z}\subset \mathbb{R}$
	 with $j$ fixed, 
	\[
	\log\left(\frac{f(n+j)}{f(n)}\right)=\log\left(\frac{g(n+j)}{g(n)}\right)+o(1)\ \ \text{as}\ \ n\to\infty.
	\]
\end{lemma}

\begin{lemma}\label{fact2}
Let $\{a_d(n)\}_{n\ge 0}$ be as in \eqref{genexample}. 
\begin{enumerate}
	\item \label{fact2part1} The sequence $\{a_d(n)\}_{n\ge 0}$ satisfies {\rm{\ref{GenGORZThm1assump1}}} for all $\a_1\le j\le \a_2$, {\rm{\ref{GenGORZThm1assump2}}}, and {\rm{\ref{genthmassump2}}}.
	\item \label{fact2part2} The sequence $\{\d(n)\}_{n\ge 0}$ satisfies {\rm{\ref{InflogconcaveThmassump1}}} for all $\a_1\le j\le \a_2$, {\rm{\ref{InflogconcaveThmassump2}}}, and {\rm{\ref{InflogconcaveThmassump3}}}. 
\end{enumerate}
\end{lemma}

\begin{proof}
%	We
%first show that $\{a_d(n)\}_{n\ge 0}$ satisfies \ref{GenGORZThm1assump1} with $0\le j\le m$ and \ref{GenGORZThm1assump2}. 
From \eqref{genexample} and \Cref{Fact1}, we have,
\begin{align}\label{neweqn9}
\log\left(\frac{a_d(n+j)}{a_d(n)}\right)
&=A(n)j-\d^2(n)j^2+\sum_{\ell\ge 3}A_{\ell}(n)j^\ell+o(1)\ \ \left(\text{ as}\ n\to \infty\right),
\end{align}
%{\bf \color{blue} K: continued from here.}	
where 
\begin{align*}
A(n)&:=\sum_{\l\in \mathcal{S}}\l A_{\l}(d)n^{\l-1}-\frac{c_2(d)}{n},\quad \d(n):=\frac1{\sqrt2n}\sqrt{\sum_{\l\in \mathcal{S}}\l(1-\l) A_{\l}(d)n^{\l}-c_2(d)},\\
A_{\ell}(n)&:=\sum_{\l\in \mathcal{S}}\binom{\l}{\ell}A_{\l}(d)n^{\l-\ell}+\frac{(-1)^\ell c_2(d)}{\ell n^\ell}.
\end{align*}

\noindent \ref{fact2part1} Using \ref{genexampleassump2} and \ref{genexampleassump3}, we have, as $n\to \infty$,  
\begin{equation}\label{obs4}
A(n)\sim \frac{\l^*A_{\l^*}(d)}{n^{1-\l^*}}\to 0.
\end{equation}
%as $n\to \infty$
%\[
%\mathcal{A}(n)\sim \frac{\l^*A_{\l^*}(d)}{n^{1-\l^*}}\to 0\ \ \text{ as }\ \ n\to \infty.
%\]	
Thus \ref{genthmassump2} is satisfied. From \ref{genexampleassump3}, we obtain $\d(n)>0$ for $n\!\gg\!1$. Using \ref{genexampleassump3}, we have 
\begin{equation}\label{impfact}
\d(n)\sim \sqrt{\frac{\l^*(1-\l^*)A_{\l^*}(d)}{2}}n^{\frac{\l^*}{2}-1}\to 0\ \ \ \ \ \ (\text{as}\ n\to \infty). %\ \ \text{ as}\ \ n\to \infty.
\end{equation}
 Thus \ref{GenGORZThm1assump2} is satisfied. Next, for  $\a_1\le j\le \a_2$, we show that
\begin{align}%\label{newlem2eqn0b}
%\d(n)&>0\ \text{and}\ \d(n)\to 0,\\%\label{newlem2eqn1}
%\mathcal{A}_\ell(n)&=o\left(\d^{\ell}(n)\right)&&\hspace{-3cm}\text{for}\ \ 3\le \ell\le 2m,\\
\label{jenseneqn1}
\sum_{\ell\ge 3}A_{\ell}(n)j^\ell&=o\left(1\right).%\ \text{with}\ \ K:=\max\{|\a_1|, |\a_2|\}.
\end{align}
Using $|\smallbinom{\l}{\ell}|\le \l$ for $\ell\in \mathbb{N}_{\ge 3}$ and \ref{genexampleassump2}\textendash\ref{genexampleassump4}, we estimate, for $\ell\in\mathbb{N}_{\ge 3}$, 
\begin{equation}\label{newlem2eqn3a}
|A_{\ell}(n)| \le \left(\l^*\sum_{\l\in \mathcal{S}} |A_{\l}(d)|+\frac{\left|c_2(d)\right|}{3}\right)n^{\l^*-\ell}=:\k_d n^{\l^*-\ell}.
\end{equation}
 Using \eqref{newlem2eqn3a}, we bound, for $n\!\gg\!1$ (and $n\ge 2K$, where $K:=\max\{|\a_1|, |\a_2|\}$), 
\begin{align*}
\left|\sum_{\ell\ge 3}A_{\ell}(n)j^\ell\right|&\le \sum_{\ell\ge 3}\left|A_{\ell}(n)\right|K^\ell
\le 2\k_{d}K^{3}n^{\l^*-3}.
\end{align*}
 Consequently, by \ref{genexampleassump2}, and \ref{genexampleassump3}, we obtain that \eqref{jenseneqn1} holds. Using this and \eqref{neweqn9}, $\{a_d(n)\}_{n\ge 0}$ satisfies \ref{GenGORZThm1assump1} with $\a_1\le j\le \a_2$. This concludes the proof of \ref{fact2part1}. 

\noindent \ref{fact2part2}  By \eqref{impfact} and \Cref{Fact1}, we have %{\bf \color{blue} K: steps in details file.}
\begin{align}\label{neweqn11}
\log\left(\frac{\d(n+j)}{\d(n)}\right)
&=B(n)j+\beta^2(n)j^2+\sum_{m\ge 3}B_{m}(n)j^m+o(1),
\end{align}
where
\[
B(n):=\left(\frac{\l^*}{2}-1\right)\frac 1n,\quad \beta(n):=\sqrt{\frac 12\left(1-\frac{\l^*}{2}\right)}\frac 1n,\quad B_m(n)=\left(\frac{\l^*}{2}-1\right)\frac{(-1)^{m+1}}{m n^m}.
\]
We have $\lim_{n\to \infty}B(n)=0$ and $\b(n)=o(\d(n))$. Thus \ref{InflogconcaveThmassump2} and \ref{InflogconcaveThmassump3} are satisfied. Note that, for $\a_1\le j\le \a_2$ and $n\!\gg\!1$, %{\bf \color{blue} K: proof in the details file.}
\[
\sum_{m\ge 3}B_{m}(n)j^m=o(1).
\] 
Using \eqref{neweqn11}, we see that $\{\d(n)\}_{n\ge 0}$ satisfies \ref{InflogconcaveThmassump1} with $\a_1\le j\le \a_2$.\qedhere 
\end{proof}

\end{document}
